\documentclass[a4paper,12pt]{article}
\usepackage{ctex}
\usepackage{geometry}
\usepackage{enumitem}
\usepackage{amsmath}
\usepackage{listings}

\geometry{left=2.5cm, right=2.5cm, top=2.5cm, bottom=2.5cm}

\title{\textbf{第二章\ 程序的控制结构}}
\author{}
\date{}

\begin{document}

\maketitle

\section*{一、选择题}

\begin{enumerate}
	\item C语言中用于结构化程序设计的三种基本结构是
	      \begin{enumerate}
		      \item[A)] 顺序结构、选择结构、循环结构
		      \item[B)] if、switch、break
		      \item[C)] for、while、do-while
		      \item[D)] if、for、continue
	      \end{enumerate}
	      \vspace{0.3cm}

	\item 以下叙述正确的是
	      \begin{enumerate}
		      \item[A)] do-while语句构成的循环不能用其它语句构成的循环来代替。
		      \item[B)] do-while语句构成的循环只能用break语句退出。
		      \item[C)] 用do-while语句构成的循环，在while后的表达式为非零时结束循环。
		      \item[D)] 用do-while语句构成的循环，在while后的表达式为零时结束循环。
	      \end{enumerate}
	      \vspace{0.3cm}

	\item 循环语句有for语句、\_\_\_\_\_\_和\_\_\_\_\_\_。
	      \begin{enumerate}
		      \item[A)] while语句
		      \item[B)] do-while语句
		      \item[C)] goto语句
		      \item[D)] 以上都是
	      \end{enumerate}
	      \vspace{0.3cm}

	\item 以下程序的输出结果是\_\_\_\_\_\_。
	      \begin{verbatim}
#include <stdio.h>
void main()
{
    int a=1,b=4,c=2;
    if(a<b)
        if(b<0) c=0;
        else c++;
    printf("%d\n",c);
}
    \end{verbatim}
	      \begin{enumerate}
		      \item[A)] 0
		      \item[B)] 1
		      \item[C)] 2
		      \item[D)] 3
	      \end{enumerate}
	      \vspace{0.3cm}

	\item 有以下程序
	      \begin{verbatim}
#include <stdio.h>
void main()
{
    int a=1,b=0;
    if(!a) b++;
    else if(a==0)
        if(a) b++;
        else b++;
    printf("%d\n",b);
}
    \end{verbatim}
	      程序运行后的输出结果是\_\_\_\_\_\_。
	      \begin{enumerate}
		      \item[A)] 0
		      \item[B)] 1
		      \item[C)] 2
		      \item[D)] 3
	      \end{enumerate}
	      \vspace{0.3cm}

	\item 以下程序的输出结果是\_\_\_\_\_\_。
	      \begin{verbatim}
#include <stdio.h>
void main()
{
    int a=0,b=1,c=0,d=20;
    if(a) d=d-10;
    else if(!b)
        if(!c) d=15;
        else d=25;
    printf("%d\n",d);
}
    \end{verbatim}
	      \begin{enumerate}
		      \item[A)] 10
		      \item[B)] 20
		      \item[C)] 30
		      \item[D)] 程序进入死循环
	      \end{enumerate}
	      \vspace{0.3cm}

	\item 下面程序的运行结果是\_\_\_\_\_\_。
	      \begin{verbatim}
#include <stdio.h>
void main()
{
    int x=1,y=0,a=0,b=0;
    switch(x)
    {
    case 1:
        switch(y)
        {
        case 0: a++; break;
        case 1: b++; break;
        }
    case 2: a++; b++; break;
    case 3: a++; b++;
    }
    printf("a=%d,b=%d\n",a,b);
}
    \end{verbatim}
	      \begin{enumerate}
		      \item[A)] a=2,b=1
		      \item[B)] a=1,b=1
		      \item[C)] a=1,b=0
		      \item[D)] a=2,b=2
	      \end{enumerate}
	      \vspace{0.3cm}

	\item 有以下程序段
	      \begin{verbatim}
int a,b,c;
a=10; b=50; c=30;
if(a>b) a=b,b=c,c=a;
printf("a=%d b=%d c=%d\n",a,b,c);
    \end{verbatim}
	      程序的输出结果是\_\_\_\_\_\_。
	      \begin{enumerate}
		      \item[A)] a=10 b=50 c=10
		      \item[B)] a=10 b=50 c=30
		      \item[C)] a=10 b=30 c=10
		      \item[D)] a=50 b=30 c=50
	      \end{enumerate}
	      \vspace{0.3cm}

	\item continue语句的作用是\_\_\_\_\_\_。
	      \begin{enumerate}
		      \item[A)] 结束整个循环
		      \item[B)] 结束本次循环，进入下一次循环
		      \item[C)] 跳出switch语句
		      \item[D)] 结束程序
	      \end{enumerate}
	      \vspace{0.3cm}

	\item break语句通常用在\_\_\_\_\_\_语句和循环语句中。
	      \begin{enumerate}
		      \item[A)] if
		      \item[B)] switch
		      \item[C)] for
		      \item[D)] while
	      \end{enumerate}
	      \vspace{0.3cm}

	\item 在C语言中，能够实现循环结构的语句有：while语句、if/goto语句、do-while语句以及\_\_\_\_\_\_语句。
	      \begin{enumerate}
		      \item[A)] for
		      \item[B)] if
		      \item[C)] switch
		      \item[D)] break
	      \end{enumerate}
	      \vspace{0.3cm}

	\item
	      \begin{verbatim}
main(){
    int n=100;
    if(n>100) printf("***");
    else printf("###");
}
    \end{verbatim}
	      输出\_\_\_\_\_\_。
	      \begin{enumerate}
		      \item[A)] \*\*\*
		      \item[B)] \#\#\#
		      \item[C)] \*\#\*
		      \item[D)] 以上都不对
	      \end{enumerate}
	      \vspace{0.3cm}

	\item
	      \begin{verbatim}
main(){
    int a=2,b=-1,c=2;
    if(a<b)
        if(b<0) c=0;
        else c+=1;
    printf("c=%d\n",c);
}
    \end{verbatim}
	      输出\_\_\_\_\_\_。
	      \begin{enumerate}
		      \item[A)] c=0
		      \item[B)] c=1
		      \item[C)] c=2
		      \item[D)] c=3
	      \end{enumerate}
	      \vspace{0.3cm}

	\item
	      \begin{verbatim}
main(){
    int k=3,n=0;
    while(k>0){
        switch(k){
            case 1: n+=k;
            case 2:
            case 3: n+=k;
            default: break;
        }
        k--;
    }
    printf("%d\n",n);
}
    \end{verbatim}
	      输出\_\_\_\_\_\_。
	      \begin{enumerate}
		      \item[A)] 0
		      \item[B)] 3
		      \item[C)] 6
		      \item[D)] 9
	      \end{enumerate}
	      \vspace{0.3cm}

	\item
	      \begin{verbatim}
main(){
    int y=10;
    while(y--);
    printf("y=%d\n",y);
}
    \end{verbatim}
	      输出\_\_\_\_\_\_。
	      \begin{enumerate}
		      \item[A)] y=10
		      \item[B)] y=1
		      \item[C)] y=随机值
		      \item[D)] y=-1
	      \end{enumerate}
	      \vspace{0.3cm}

	\item
	      \begin{verbatim}
#include<stdio.h>
main(){
    int s=0,m;
    for(m=7;m>=3;m--)
        switch(m){
            case 1: case 4: case 7:s++;break;
            case 2: case 3: case 6:s+=2;
            case 0: case 5:s+=3;break;
        }
    printf("s=%d\n",s);
}
    \end{verbatim}
	      输出\_\_\_\_\_\_。
	      \begin{enumerate}
		      \item[A)] s=0
		      \item[B)] s=3
		      \item[C)] s=6
		      \item[D)] s=9
	      \end{enumerate}
	      \vspace{0.3cm}

	\item 设 \texttt{int a=12}，则执行完语句 \texttt{a+=a-=a*a;} 后a的值是\_\_\_\_\_\_。
	      \begin{enumerate}
		      \item[A)] 552
		      \item[B)] 264
		      \item[C)] 144
		      \item[D)] -264
	      \end{enumerate}
	      \vspace{0.3cm}

	\item 以下程序的输出结果是\_\_\_\_\_\_。
	      \begin{verbatim}
#include <stdio.h>
void main()
{
    int a=5,b=4,c=3,d;
    d=(a>b>c);
    printf("%d\n",d);
}
    \end{verbatim}
	      \begin{enumerate}
		      \item[A)] 0
		      \item[B)] 1
		      \item[C)] 非0的数
		      \item[D)] -1
	      \end{enumerate}
	      \vspace{0.3cm}

	\item 为表示关系 $x \ge y \ge z$，应使用的C语言表达式是\_\_\_\_\_\_。
	      \begin{enumerate}
		      \item[A)] (x>=y)\&\&(y>=z)
		      \item[B)] (x>=y) AND(y>=z)
		      \item[C)] (x>=y>=z)
		      \item[D)] (x>=y)\&(y>=z)
	      \end{enumerate}
	      \vspace{0.3cm}

	\item 若有以下程序段：\\
	      \texttt{int a[4]={0,0,0,0}, i;}\\
	      \texttt{for(i=1; i<4; i++) scanf("\%d", \&a[i]);}\\
	      若从键盘输入：1234<回车>，则 \texttt{a[1]} 的值是\_\_\_\_\_\_。
	      \begin{enumerate}
		      \item[A)] 不定值
		      \item[B)] 0
		      \item[C)] 1234
		      \item[D)] 4
	      \end{enumerate}
	      \vspace{0.3cm}

	\item C语言中，break语句通常用在\_\_\_\_\_\_语句和循环语句中。
	      \begin{enumerate}
		      \item[A)] if
		      \item[B)] switch
		      \item[C)] for
		      \item[D)] B和C
	      \end{enumerate}
	      \vspace{0.3cm}

	\item 若有定义 \texttt{a[10];} 则允许数组a的下标最小可以是\_\_\_\_\_\_。
	      \begin{enumerate}
		      \item[A)] 0
		      \item[B)] 1
		      \item[C)] -1
		      \item[D)] 任意整数
	      \end{enumerate}
	      \vspace{0.3cm}

	\item 若有 \texttt{a=3, b=5;} 则求 \texttt{a>b} 的关系运算结果是\_\_\_\_\_\_。
	      \begin{enumerate}
		      \item[A)] 0
		      \item[B)] 1
		      \item[C)] 真
		      \item[D)] 假
	      \end{enumerate}
	      \vspace{0.3cm}
\end{enumerate}

\section*{二、填空题}

\begin{enumerate}
	\item 打印出100到200之间的所有素数（即除1和它本身再没有其他约数的数）。
	      \begin{verbatim}
#include <stdio.h>
#include <math.h>
int main() {
    int i, j;
    for (i = 100; i <= 200; i++) {
        int temp = (int)sqrt(i);
        for (j = 2; j <= temp; j++)
            if (i % j == 0) break;
        if (j > temp)
            printf("%d ", i);
    }
    printf("\n");
    return 0;
}
    \end{verbatim}
	      \vspace{0.3cm}

	\item 斐波那契数列的第1和第2个数分别为0和1，从第三个数开始，每个数等于其前两个数之和。求斐波那契数列中的前20个数，要求每行输出5个数。
	      \begin{verbatim}
#include <stdio.h>
int main() {
    int f, f1, f2, i;
    printf("斐波那契数列：\n");
    f1 = 0; f2 = 1;
    printf("%6d%6d", f1, f2);
    for (i = 3; i <= 20; i++) {
        f = f1 + f2;
        printf("%6d", f);
        if (i % 5 == 0) printf("\n");
        f1 = f2;
        f2 = f;
    }
    printf("\n");
    return 0;
}
    \end{verbatim}
	      \vspace{0.3cm}

	\item 计算 $1 - \frac{1}{2} + \frac{1}{3} - \frac{1}{4} + \cdots - \frac{1}{10}$ 的值。
	      \begin{verbatim}
#include <stdio.h>
int main() {
    double x, p1 = 1, p2 = 1, s = 0;
    int i, j = 1;
    printf("输入 x 的值:");
    scanf("%lf", &x);
    for (i = 1; i <= 10; i++) {
        p1 *= x;
        p2 *= i;
        s += j * p1 / p2;
        j = -j;
    }
    printf("%f\n", s);
    return 0;
}
    \end{verbatim}
	      \vspace{0.3cm}

	\item 采用辗转相除法求出两个整数的最大公约数。
	      \begin{verbatim}
#include <stdio.h>
int main() {
    int a, b;
    printf("请输入两个正整数:");
    scanf("%d%d", &a, &b);
    while (a <= 0 || b <= 0) {
        printf("重新输入:");
        scanf("%d%d", &a, &b);
    }
    while (b) {
        int r = a % b;
        a = b;
        b = r;
    }
    printf("最大公约数：%d\n", a);
    return 0;
}
    \end{verbatim}
	      \vspace{0.3cm}

	\item 把从键盘上输入的一个大于等于3的整数分解为质因子的乘积。
	      \begin{verbatim}
#include <stdio.h>
int main() {
    int x;
    printf("请输入一个整数，若小于3则重输:");
    do {
        scanf("%d", &x);
    } while (x < 3);
    int i = 2;
    do {
        while (x % i == 0) {
            printf("%d ", i);
            x /= i;
        }
        i++;
    } while (i < x);
    if (x != 1) printf("%d", x);
    printf("\n");
    return 0;
}
    \end{verbatim}
	      \vspace{0.3cm}

	\item 在输出屏幕上打印出一个由字符'*'组成的等腰三角形，该三角形的高为5行，从上到下每行的字符数依次为1,3,5,7,9。
	      \begin{verbatim}
#include <stdio.h>
int main() {
    int i, j;
    for (i = 1; i <= 5; i++) {
        for (j = 1; j <= 9; j++)
            if (j <= 5 - i || j >= 5 + i)
                printf(" ");
            else
                printf("*");
        printf("\n");
    }
    return 0;
}
    \end{verbatim}
	      \vspace{0.3cm}

	\item 求1!+2!+...+20!的值。
	      \begin{verbatim}
#include <stdio.h>
int main() {
    double s = 0, t = 1;
    int n;
    for (n = 1; n <= 20; n++) {
        t = t * n;
        s = s + t;
    }
    printf("1!+2!+...+20! = %e\n", s);
    return 0;
}
    \end{verbatim}
	      \vspace{0.3cm}

	\item 一球从100m高度自由落下，每次落地后反跳回原高度的一半，再落下。求它在第10次落地时，共经过多少m？第10次反弹多高？
	      \begin{verbatim}
#include <stdio.h>
int main() {
    float sn = 100, hn = 100;
    int n;
    for (n = 2; n <= 10; n++) {
        sn = sn + hn;
        hn = hn / 2;
    }
    printf("第10次落地时共经过%f m.\n", sn);
    printf("第10次反弹%f m.\n", hn);
    return 0;
}
    \end{verbatim}
	      \vspace{0.3cm}

	\item 下面程序的运行结果是\_\_\_\_\_\_。
	      \begin{verbatim}
main()
{ int i,s;
  for(i=10,s=0;i;s+=i,i--)
      printf("result:%d\n",s);
}
    \end{verbatim}
	      \vspace{0.3cm}

	\item
	      \begin{verbatim}
main(){
    int i,j,x;
    for(i=0;i<=4;i++){
        for(j=1;j<=4-i;j++) printf(" ");
        for(j=0;j<=2*i+1;j++) printf("*");
        printf("\n");
    }
}
    \end{verbatim}
	      输出\_\_\_\_\_\_。
	      \vspace{0.3cm}

	\item 下面程序的功能是找出100到200之间不能被3整除但能被5整除的数。
	      \begin{verbatim}
#include <stdio.h>
int main() {
    int m;
    for (m = 100; m <= 200; m++)
        if (m % 3 != 0 && m % 5 == 0)
            printf("%d\t", m);
    printf("\n");
    return 0;
}
    \end{verbatim}
	      \vspace{0.3cm}

	\item 输入a、b、c三个值，输出其中最大者。
	      \begin{verbatim}
#include <stdio.h>
int main() {
    int a, b, c, max;
    printf("请输入三个数 a,b,c:\n");
    scanf("%d,%d,%d", &a, &b, &c);
    max = a;
    if (max < b)
        max = b;
    if (max < c)
        max = c;
    printf("最大数为：%d", max);
    return 0;
}
    \end{verbatim}
	      \vspace{0.3cm}

	\item 统计字符串中英文字母个数的程序。
	      \begin{verbatim}
#include <stdio.h>
int count(char str[]) {
    int num = 0;
    for (int i = 0; str[i]; i++)
        if ((str[i] >= 'a' && str[i] <= 'z') || (str[i] >= 'A' && str[i] <= 'Z'))
            num++;
    return num;
}
int main() {
    char s1[80];
    printf("Enter a line:");
    scanf("%s", s1);
    printf("count=%d\n", count(s1));
    return 0;
}
    \end{verbatim}
	      \vspace{0.3cm}

	\item 本函数用对分查找法，在已按字母次序从小到大排序的字符数组list中查找字符c，若c在数组中，函数返回字符c在数组中的下标，否则返回-1。
	      \begin{verbatim}
int search(char list[], char c, int len) {
    int low, high, k;
    low = 0; high = len - 1;
    while (low <= high) {
        k = (low + high) / 2;
        if (list[k] == c) return k;
        else if (list[k] > c) high = k - 1;
        else low = k + 1;
    }
    return -1;
}
    \end{verbatim}
	      \vspace{0.3cm}

	\item 函数mycmp(char *s, char *t)的功能是比较字符串s和t的大小，当s等于t时返回0，否则返回s和t的第一个不同字符的ASCII码的差值。
	      \begin{verbatim}
int mycmp(char *s, char *t) {
    while (*s == *t) {
        if (*s == '\0') return 0;
        ++s; ++t;
    }
    return *s - *t;
}
    \end{verbatim}
	      \vspace{0.3cm}

	\item 指出程序的错误并改正：
	      \begin{verbatim}
#include <stdio.h>
int main() {
    int a, b, max;
    scanf("%d,%d", &a, &b);
    if (a < b) max = a;
    max = b;
    printf("max=%d", max);
    return 0;
}
    \end{verbatim}
	      \vspace{0.3cm}

	\item 下面程序能够完成交换数组a和数组b中的对应元素的功能。
	      \begin{verbatim}
#include <stdio.h>
void swap(int *p1, int *p2) {
    int temp;
    temp = *p1;
    *p1 = *p2;
    *p2 = temp;
}
int main() {
    int a[5] = {1, 3, 5, 7, 9};
    int b[5] = {2, 4, 6, 8, 10};
    int i;
    for (i = 0; i < 5; i++) swap(&a[i], &b[i]);
    for (i = 0; i < 5; i++) printf("a[%d]=%-4d", i, a[i]);
    printf("\n");
    for (i = 0; i < 5; i++) printf("b[%d]=%-4d", i, b[i]);
    printf("\n");
    return 0;
}
    \end{verbatim}
	      \vspace{0.3cm}

	\item 某大学举行的演讲比赛中，有十个评委为参赛的选手打分，分数为0～100分。选手最后得分为：去掉一个最高分和一个最低分后其余八个分数的平均值。
	      \begin{verbatim}
#include <stdio.h>
int main() {
    int max, min, score, i;
    int sum = 0;
    max = 0; min = 100;
    for (i = 0; i < 10; i++) {
        printf("Input number %d=", i + 1);
        scanf("%d", &score);
        sum += score;
        if (max < score) max = score;
        if (min > score) min = score;
        printf("\n");
    }
    printf("Canceled max score:%d\n", max);
    printf("Canceled min score:%d\n", min);
    printf("average score:%lf\n", (sum - max - min) / 8.0);
    return 0;
}
    \end{verbatim}
	      \vspace{0.3cm}

	\item 把指针str所指的字符串按相反的顺序赋给rev[]。
	      \begin{verbatim}
#include <stdio.h>
#include <string.h>
int main() {
    char *str = "abcdefg";
    char rev[10];
    int i;
    for (i = 0; i < 7; i++)
        rev[i] = str[6 - i];
    rev[i] = '\0';
    printf("%s\n", rev);
    return 0;
}
    \end{verbatim}
	      \vspace{0.3cm}
\end{enumerate}

\section*{三、判断题}

\begin{enumerate}
	\item C语言程序总是从源程序文件中的第一个函数开始执行。( \ \ )
	      \vspace{0.3cm}

	\item 数组名代表数组所占存储区的首地址，其值不可以改变。( \ \ )
	      \vspace{0.3cm}

	\item elseif不属于C语言关键字（保留字）( \ \ )
	      \vspace{0.3cm}

	\item 指针变量可以加减一个整数。( \ \ )
	      \vspace{0.3cm}

	\item 宏替换不占用运行时间。( \ \ )
	      \vspace{0.3cm}

	\item C语言中转义字符以"$\setminus\setminus$"开头。( \ \ )
	      \vspace{0.3cm}

	\item C语言规定，函数返回值的类型是由return语句中的表达式类型决定的。( \ \ )
	      \vspace{0.3cm}

	\item 如果在一个函数中的复合语句中定义了一个变量，则该变量只在该复合语句中有效。( \ \ )
	      \vspace{0.3cm}

	\item C语言中的函数既可以递归定义，又可以嵌套定义。( \ \ )
	      \vspace{0.3cm}

	\item main函数可以有参数。( \ \ )
	      \vspace{0.3cm}
\end{enumerate}

\section*{四、程序填空题}

\begin{enumerate}
	\item 本程序能够在屏幕中央显示出如下图形。
	      \begin{verbatim}
#######
#####
###
#
    \end{verbatim}
	      填空使程序完整：
	      \begin{verbatim}
#include <stdio.h>
int main() {
    int i, j, k;
    for (i = 4; i >= 1; i--) {
        for (k = 1; k <= 4 - i; k++)
            printf(" ");
        for (j = 1; j <= 2 * i - 1; j++)
            printf("#");
        printf("\n");
    }
    return 0;
}
    \end{verbatim}
	      \vspace{0.3cm}

	\item 下面程序通过指向整型变量的指针将数组m[4][3]的内容按4行3列的格式输出，请填空使程序完整。
	      \begin{verbatim}
#include <stdio.h>
int main() {
    int m[4][3] = {{1,2,3},{4,5,6},{7,8,9},{10,11,12}};
    int i, j, *p = m[0];
    for (i = 0; i < 4; i++) {
        for (j = 0; j < 3; j++) printf("%4d", *(p + i * 3 + j));
        printf("\n");
    }
    return 0;
}
    \end{verbatim}
	      \vspace{0.3cm}
\end{enumerate}

\end{document}
